Recall that the $s$-\textbf{th cyclotomic polynomial} $\Phi_s (X)$ is the unique monic, irreducible polynomial whose roots are the primitive $s$-th roots of unity.
In other words, $z\in\mathbb{C}$ is a root of $\Phi_s(z)=0$ if and only if we have $z=e^{2\pi i d/s}$ for some $d\in\ZZ$ with $(d,s)=1$. 
Alternatively, $\Phi_s$ can be defined inductively via the identity
\begin{equation}\label{poly-e0}
X^N-1=\prod_{s|N}\Phi_s(X).
\end{equation}
Thus $\Phi_1(X)=X-1$, and 
\begin{equation}\label{polysec-e0}
1+X+X^2+\dots+X^{N-1}=\prod_{s|N,s\neq 1}\Phi_s(X)\hbox{ for }N>1.
\end{equation} 
In particular, if $p$ is a prime number, then $\Phi_p(X)=1+X+\dots+X^{p-1}$.
